\subsection{标准群†}

若 \(\mathrm{S}(\mathbf{X}_{1},\mathbf{X}_{2},\ldots,\mathbf{X}_{n})\) 是系数在 A 中的形式幂级数，则对于 m 中的所有 \(x_{1},\ldots,x_{r}\)，级数 \(\mathrm{S}(x_{1},x_{2},\ldots,x_{r})\) 收敛。更确切地说，\(m\times m\times\cdots\times m\) 包含于 S 的绝对收敛域（可微与解析流形，R，4.1.3）。

定义 1. 设 \(r\) 是满足 \(\geqslant 0\)\(r\)\(\mathbf{K}\)\(G\) 的整数。K 上 r 维的标准群是具有下列性质的李群 G：

\begin{enumerate}
\def\labelenumi{(\roman{enumi})}
\item
  G 的底层解析流形是 m × m × · · · × m（r 个因子）；
\item
  存在一个以 \(\mathbf{F}\) 为系数、无常数项的关于 \(2r\) 个变量的形式幂级数 \(\mathbf{A}^r\)，使得对 G 中所有 \(x.y = \mathbf{F}(x,y)\)，都有 \(x,y\)\(\mathbf{G}\)。
\end{enumerate}

于是 0.0 = 0，从而 G 的单位元是 m × m · · · × m 的原点。
† 第 3、4 号的结果及其证明在 K 的特征为正时仍然成立。

将 L(G) 识别为 \(K^r\)。由 § 5，公式 (13)，L(G) 关于典范基的结构常数属于 A。在同一证明中，我们需要把 m × ⋯ × m 中的元素既看作 G 的元素，又看作 L(G) 的元素。

例。设 \(G = 1 + \mathbf{M}_{n}(m)\)，它是 \(\mathbf{M}_{n}(\mathbf{K})\) 的开子集。若 \(x \in G\)，则 det \(x \in 1 + m\)，因而 \(G \subset \mathbf{GL}(n, K)\)。显然 \(GG \subset G\)。若 \(x = 1 + y\)，其中 \(y \in \mathbf{M}_{n}(m)\)，则矩阵求逆的计算首先说明 \(x^{-1} \in \mathbf{M}_{n}(\mathbf{A})\)；若写作 \(x^{-1} = 1 + y'\)，则 \(y + y' + yy' = 0\)，从而 \(y' \in \mathbf{M}_{n}(m)\)，因此 \(x^{-1} \in G\)。所以 G 是 \(\mathbf{GL}(n, K)\) 的开子群。我们借助映射 \(m^{n^2}\) 将 G 与 \((\delta_{ij} + y_{ij}) \mapsto (y_{ij})\) 识别。显然 G 是标准群。

定理 4. 设 G 是有限维李群。G 存在一个开子群同构于标准群。

将 G 替换为与 G 的某个开子群同构的群后，问题归结为 G 是 \(K^{r}\) 的开子集，单位元为 0，且乘积 x.y 和逆元 \(x^{[-1]}\) 的坐标由下列公式给出：

\begin{enumerate}
\def\labelenumi{(\arabic{enumi})}
\setcounter{enumi}{2}
\tightlist
\item
\end{enumerate}

\[
(x. y) _ {i} = x _ {i} + y _ {i} + \sum_ {| \alpha | \geqslant 1, | \beta | \geqslant 1} c _ {\alpha \beta i} x ^ {\alpha} y ^ {\beta} (i = 1, 2, \dots , r)\tag{3}
\]

\[
(x ^ {[ - 1 ]}) _ {i} = - x _ {i} + \sum_ {| \alpha | > 1} d _ {\alpha i} x ^ {\alpha} \quad (i = 1, 2, \dots , r)\tag{4}
\]

令 \(x, y\)\(G\)\(\lambda \in K^*\)\(G\)，通过比为 \(G' = \lambda G\) 的位似将群运算从 G 搬到 \(\lambda\)。对于 G' 中的 \(x', y'\)\(G'\)，在 G' 中计算的乘积 \(x'.y'\) 和逆元 \(x'^{(-1)}\)\(G'\) 的坐标为

\[
(x ^ {\prime}. y ^ {\prime}) _ {i} = x _ {i} ^ {\prime} + y _ {i} ^ {\prime} + \sum_ {| \alpha | \geqslant 1, | \beta | \geqslant 1} c _ {\alpha \beta i} ^ {\prime} x ^ {\prime \alpha} y ^ {\prime \beta} (i = 1, 2, \dots , r)
\]

\[
(x ^ {\prime [ - 1 ]}) _ {i} = - x _ {i} ^ {\prime} + \sum_ {| \alpha | > 1} d _ {\alpha i} ^ {\prime} x ^ {\prime \alpha} \quad (i = 1, 2, \dots , r)
\]

其中

\[
c _ {\alpha \beta i} ^ {\prime} = \lambda^ {- | \alpha | - | \beta | + 1} c _ {\alpha \beta i}, \quad d _ {\alpha i} ^ {\prime} = \lambda^ {- | \alpha | + 1} d _ {\alpha i}.
\]

由于级数 (3) 和 (4) 收敛可知，当 \(|\lambda|\) 充分大时，对所有 \(\alpha, \beta, i\)，

\[
\left| c _ {\alpha \beta i} ^ {\prime} \right| \leqslant 1, \quad \left| d _ {\alpha i} ^ {\prime} \right| \leqslant 1
\]

即 \(c_{\alpha\beta i}^{\prime}\in A\) 且 \(d_{\alpha i}^{\prime}\in A\)；另一方面 \(G^{\prime}\supset m\times m\times\cdots\times m\)。于是 \(m\times m\times\cdots\times m\) 是 \(G^{\prime}\) 的开子群，且是标准群。

